% !TEX root = ../main.tex
\chapter{Сурах бичгийн шаардлага}\label{ch:req}

\section{Номын үндсэн шаардлага}
Бүхэл тоон жагсаалтыг ижил хэмжээтэй хоёр хэсэгт хувааж, хоёр тусдаа thread тус бүр нэг хэсгийг дурын алгоритмаар
эрэмбэлнэ. Гурав дахь thread хоёр хэсгийг эх массивтай ижил хэмжээтэй хоёр дахь global массив руу нэгтгэнэ.
Parent thread бүх sorting thread дуусахыг хүлээж, дараа нь merging thread-ийг эхлүүлж, түүнийг дуустал хүлээгээд
үр дүнг хэвлэнэ. Thread бүрд эрэмбэлэх хэсгийн эхлэх индекс зэрэг параметрийг дамжуулна~\cite[Ch.~4, Programming Projects]{osc10}.

\section{Нэмэлт шаардлага}
Сэдвийн хуваарилалтаар \emph{Sorting animation} нэмэгдсэн: дүрслэл нь бодит thread-тэй холбоотой байж, хэмжээ ба
утгыг өөрчлөх, Run / Step / Pause / Reset, хурдны удирдлагатай байна. Мөн багшийн заавраар гарын авлага, тайлан,
GitHub repository шаардлагатай. Баг өөрөө оролтын бүрэн шалгалт (хоосон, тоо биш, хэт том утга, $n = 0\dots3$,
давхардал, сондгой $n$) болон зөвхөн стандарт сан ашиглах шаардлагыг тавьсан.

\section{Шаардлагын хангалт}
\begin{table}[H]
\centering\small
\caption{Номын шаардлага ба хэрэгжүүлэлт}\label{tab:req}
\begin{tabularx}{\linewidth}{@{}YY@{}}
\toprule
Шаардлага & Хэрэгжүүлэлт \\
\midrule
Жагсаалтыг хоёр хуваах & \code{splitIndex(n)} $=\lceil n/2\rceil$; зүүн \code{[0, mid)}, баруун \code{[mid, n)} \\
Хоёр sorting thread & \code{SortingThread extends Thread}; 5 алгоритмаас хэсэг тус бүрд сонгоно \\
Global жагсаалт ба хоёр дахь массив & \code{SharedData.list}, \code{SharedData.result = new int[n]} \\
Merging thread & \code{MergingThread}: хоёр заагчтай, тогтвортой нэгтгэл \\
Parent хүлээгээд нэгтгэнэ & \code{left.join(); right.join(); merger.start(); merger.join();} \\
Параметр дамжуулах & \code{SortParams(name, startIndex, endIndex, algorithm)}, \code{MergeParams} \\
Үр дүн хэвлэх & Үр дүн, thread бүрийн статистик, \code{Arrays.sort}-той харьцуулсан шалгалт \\
\bottomrule
\end{tabularx}
\end{table}

\section{Номын жишээ ба сондгой $n$}
Номын жишээ \code{7, 12, 19, 3, 18, 4, 2, 6, 15, 8}-ийн хоёр хэсэг \code{3 7 12 18 19 | 2 4 6 8 15} болж эрэмбэлэгдээд,
\code{2 3 4 6 7 8 12 15 18 19} гэж нэгтгэгдэнэ. Номд хэсгүүдийг “ижил хэмжээтэй” гэсэн ч сондгой $n$ үед энэ
боломжгүй. Бид сондгой $n$-ийг алдаа гэж үзэхийн оронд хэсгүүдийг $\lceil n/2\rceil$ ба $\lfloor n/2\rfloor$ болгосон.
Нэгтгэх алгоритм тэнцүү хэмжээ шаарддаггүй бөгөөд ялгаа нь ердөө 1 элемент юм. $n = 1$ үед баруун хэсэг хоосон
байх ч номын бүтцийг хадгалж Sorter-R thread ажиллаад шууд дуусна.
